\section{非楕円成分の一様な量子化と周囲の相対対角}
\label{nq:sec:quantization}
本節では前節の固定した二つのPoisson成分を、実在する二パラメータの代数族として量子化する。
形式級数の存在や有限次数のHilbert関数検査だけにとどめず、全次数のflatnessと
\emph{周囲の代数}の有限相対対角分解まで構成する。
残余dg圏の対角は別の対象であり、本節の結論をそのsaturationと読み替えない。

\subsection{一次資料との関係と主定理}
三変数の同次Jacobian代数のAS-regular性はMori--Smith
\cite[Theorem 4.2, Lemma 4.1]{MSNE12}の分類を用いる。
中央一次元による拡大にはLe Bruyn--Smith--Van den Bergh
\cite[Theorem 2.6, Corollary 2.7, Theorem 3.1.3]{LBSVNE12}、
低次ポテンシャルを加えるPBW変形とその双加群複体にはBerger--Taillefer
\cite[Theorems 3.1 and 3.6]{BTNE12}を用いる。
これらの一般定理を新規な成果とは主張しない。
本節では、前節の特定の節点・尖点成分に対して中心混合項を保った具体的な式を与え、
係数環、例外集合、底変更および対角の種類を明記する。
中心三次元の表示の初出は未確定である。

$r,s$を次数零、$x,y,z,a,b,c,u$を次数一とし、
\begin{equation}
 R_{\mathrm C}=\CC[r,s],\qquad
 R_{\mathrm N}=\CC[r,s,(4-r^2)^{-1}]
\end{equation}
と置く。第一ブロックの代数$D_\epsilon$は$u$を中心として
\begin{align}
 \epsilon=\mathrm C:\quad
 [y,z]&=r(z^2+2ux),&[z,x]&=3ry^2,&[x,y]&=r(xz+zx),
 \label{nq:eq:cusp}\\
 \epsilon=\mathrm N:\quad
 [y,z]&=\tfrac r2(yz+zy)+2rux,&
 [z,x]&=\tfrac r2(zx+xz)+3ry^2,&
 [x,y]&=\tfrac r2(xy+yx)+3rz^2
 \label{nq:eq:node}
\end{align}
で定める。第二ブロック$E_s$も$u$を中心とし、
\begin{equation}\label{nq:eq:fermat}
 [a,b]=s(3c^2+2uc),\quad
 [b,c]=s(3a^2+2ua),\quad
 [c,a]=s(3b^2+2ub)
\end{equation}
で定める。異なるブロックの生成元は可換とする。したがって
\begin{equation}
 \mathcal A_\epsilon=D_\epsilon\otimes_{R_\epsilon[u]}E_s
\end{equation}
は七生成元の次数付き代数である。

\begin{theorem}[二つの非楕円成分の代数的量子化]
\label{nq:thm:family}
以下の式で定める$\Omega_\epsilon\in\mathcal A_{\epsilon,3}$は中心的であり、
$\mathcal B_\epsilon=\mathcal A_\epsilon/(\Omega_\epsilon)$は次を満たす。
\begin{enumerate}
\item $\mathcal A_\epsilon$は次数ごとに有限局所自由で、その各幾何学的ファイバーは
noetherianなKoszul AS-regular domain、global dimension七、AS parameter七である。
\item $\Omega_\epsilon$は各幾何学的ファイバーで非零の中心正則元であり、
$\mathcal B_\epsilon$は次数ごとに有限局所自由で、形成は任意の底変更と両立する。
各ファイバーのHilbert級数とAS-Gorenstein型は
\begin{equation}
 H_{\mathcal B}(t)=\frac{1-t^3}{(1-t)^7},\qquad
 \operatorname{injdim}\mathcal B=6,\qquad \text{AS parameter}=4.
\end{equation}
\item $(r,s)=(0,0)$で$\mathcal B_\epsilon$は前節の$X_\epsilon$の斉次座標環となる。
任意の$(\lambda,\mu)$に沿う$r=h\lambda,s=h\mu$への引戻しは、
$\lambda\pi_{P_1}+\mu\pi_{P_2}$を半古典極限に持つ。
\item $R_\epsilon[u]$上の$\mathcal A_\epsilon$の対角は、階数列
$1,6,15,20,15,6,1$の有限自由双加群分解を持つ。
この分解は任意の可換$R_\epsilon[u]$-代数への底変更後も完全である。
$R_\epsilon$上でも周囲の対角はperfectである。
\end{enumerate}
この定理は$\operatorname{qgr}\mathcal B_\epsilon$やその残余圏のsmoothnessを含まない。
\end{theorem}

\subsection{第一ブロックのmonic PBW表示}
この小節では$r=h$と書き、係数環を
$Q_{\mathrm C}=\CC[h]$、$Q_{\mathrm N}=\CC[h,(4-h^2)^{-1}]$とする。
$u$を係数に入れ、active語の長さ、次いで辞書式$x>y>z$の順序を用いる。
$\epsilon=\mathrm C$の関係式は
\begin{align}
 yz&\longmapsto zy+hz^2+2hux,\notag\\
 xz&\longmapsto zx-3hy^2,\notag\\
 xy&\longmapsto yx+2hzx-3h^2y^2
 \label{nq:eq:cusp-rewrite}
\end{align}
と等価である。節点型では$q=(2+h)/(2-h)$として
\begin{align}
 yz&\longmapsto qzy+\frac{4h}{2-h}ux,\notag\\
 xz&\longmapsto q^{-1}zx-\frac{6h}{2+h}y^2,\notag\\
 xy&\longmapsto qyx+\frac{6h}{2-h}z^2.
 \label{nq:eq:node-rewrite}
\end{align}
全ての右辺のactive語は左辺より小さい。重なりの曖昧さは$xyz$の一つだけである。
\begin{lemma}[全次数と任意の底変更に対するPBW性]
\label{nq:lem:pbw}
二つの還元系はいずれも合流的であり、
$\{z^iy^jx^k:i,j,k\ge0\}$は$D_\epsilon$の$Q_\epsilon[u]$-基底である。
この主張は任意の可換$Q_\epsilon[u]$-代数への底変更後も成立する。
\end{lemma}
\begin{proof}
節点型で二つの還元経路から得る共通の正規形は
\begin{equation}
 qzyx+\frac{4h}{2-h}ux^2+\frac{6h}{2-h}(z^3-y^3).
\end{equation}
尖点型では両経路が
\begin{equation}\begin{aligned}
 zyx+3hz^2x+2hux^2-3hy^3-9h^2zy^2-6h^3z^2y
 -6h^4z^3\\
 {}-12h^3uyx-24h^4uzx+36h^5uy^2
\end{aligned}\end{equation}
となる。よってDiamond lemma \cite[Theorem 1.2]{BergNE12}を適用できる。
係数は表示した環に属し、還元の先頭係数は全て一である。
同じ有限恒等式が任意の底変更後にも成り立つため、基底の主張も保存される。
\end{proof}

\begin{lemma}[幾何学的ファイバーの正則性]\label{nq:lem:regular4}
$D_\epsilon$の各幾何学的ファイバーは四次元noetherian Koszul AS-regular domainである。
\end{lemma}
\begin{proof}
$u=0$の三変数代数は、尖点型では全$h$、節点型では$h\ne\pm2$で
三次元Calabi--Yau AS-regularである\cite[Lemma 4.1, Theorem 4.2]{MSNE12}。
域性とnoetherian性は、次のOre表示からも確認できる。
節点型では量子平面$yz=qzy$上に
\begin{equation}
 \sigma(y)=qy,\quad\sigma(z)=q^{-1}z,\quad
 \delta(y)=\frac{6h}{2-h}z^2,\quad
 \delta(z)=-\frac{6h}{2+h}y^2
\end{equation}
による$x$のOre拡大を取る。
尖点型ではJordan平面$yz-zy=hz^2$上に
\begin{equation}
 \sigma(y)=y+2hz,\quad\sigma(z)=z,\quad
 \delta(y)=-3h^2y^2,\quad\delta(z)=-3hy^2
\end{equation}
を使う。いずれも$\sigma$は可逆、$\delta$は表示した関係を保つ$\sigma$-導分である。
補題\ref{nq:lem:pbw}より$u$は中心正則元である。
中央次数一拡大の定理\cite[Theorem 2.6, Corollary 2.7 and Remark 2.8]{LBSVNE12}
により四次元の結論を得る。ここで特異なのは三次曲線であって、周囲のAS代数ではない。
\end{proof}

\subsection{混合項を含む中心三次元の閉公式}
全ての積を、表示した順序の非可換積として読む。
\begin{proposition}[第一ブロックの中心元]\label{nq:prop:center}
次の元は$D_\epsilon$の中心に属し、$h=0$で$f_\epsilon+ux^2$となる。
\begin{align}
 \Gamma_{\mathrm C}(h)={}&
 z^2x+y^3-3hzy^2+ux^2
 -6h^2uyx-6h^3uzx+12h^4uy^2,
 \label{nq:eq:centerC}\\
 \Gamma_{\mathrm N}(h)={}&
 z^3+\left(\frac{2-h}{2+h}\right)^3y^3
 +\frac{h^2+12}{3(2+h)^2}zyx
 +\frac{h^2+12}{6(2+h)}ux^2\notag\\
 &-\frac{24h^2}{(2+h)^2}uzy
 -\frac{24h^3}{(2+h)^2}u^2x.
 \label{nq:eq:centerN}
\end{align}
\end{proposition}
\begin{proof}
補題\ref{nq:lem:pbw}の還元を両辺の積$\Gamma_\epsilon a$、$a\Gamma_\epsilon$
（$a=x,y,z$）へ適用すると、全てで同じ正規形となる。
付属の\texttt{nonelliptic\_quantization\_certificate.json}に、六つの共通正規形を全係数で収録した。
その係数は尖点で$\mathbb Q[h,u]$、節点で$\mathbb Q[h,u,(4-h^2)^{-1}]$に属し、
コードは差の分子が恒等的に零であることを検査する。
従ってこれは有限個の特殊化による観察ではなく、合流する還元による普遍恒等式の証明である。
$u$との可換性は定義による。$h=0$での式は直接代入で分かる。
\end{proof}
\begin{remark}[量子的な混合補正の必要性]
$\Gamma_\epsilon(h)|_{u=0}+ux^2$という素朴な持上げは一般には中心でない。
付属コードはその非零な交換子の正規形も記録する。
中心混合項を古典式のまま加えるだけでは、量子化の中心三次方程式にならない。
\end{remark}

第二ブロックでは
$\alpha=a+u/3$、$\beta=b+u/3$、$\gamma=c+u/3$、$t=3s$と置く。
関係式\eqref{nq:eq:fermat}は
$[\alpha,\beta]=t(\gamma^2-u^2/9)$とその巡回置換になる。
\begin{lemma}[Fermat側の混合中心元]\label{nq:lem:fermat-center}
\begin{equation}
 \Gamma_{\mathrm F}(s)=
 \alpha^3+\beta^3+\gamma^3
 -3s^2\!\sum_{w\in\mathfrak S_3(\alpha\beta\gamma)}w
 -\frac{u^2}{3}(\alpha+\beta+\gamma)+\frac{2u^3}{9}
 \label{nq:eq:centerF}
\end{equation}
は$E_s$の中心元であり、$s=0$で$a^3+b^3+c^3+u(a^2+b^2+c^2)$となる。
\end{lemma}
\begin{proof}
独立な中心係数$\eta$を用い、
$r_1=xy-yx-t(z^2-\eta)$、$r_2=yz-zy-t(x^2-\eta)$、
$r_3=zx-xz-t(y^2-\eta)$とする。
$D=x^3+y^3+z^3$、$T=xyz+xzy+yxz+yzx+zxy+zyx$に対し、
自由代数$\mathbb Z[t,\eta]\langle x,y,z\rangle$で
\begin{equation}\begin{aligned}
 [3D-t^2T-9\eta(x+y+z),x]={}&
 -3r_1y^2+t^2r_1zx-t^2xr_1z-3yr_1y-t^2zr_1x\\
 &+t^2xzr_1-3y^2r_1+2t^2r_2x^2-3tr_2yz\\
 &-t^2xr_2x-t^2x^2r_2+3tyzr_2
 +3r_3z^2+3zr_3z+3z^2r_3\\
 &+9\eta r_1-9\eta r_3
\end{aligned}\end{equation}
が成立する。語ごとの展開と巡回置換により三交換子を処理できる。
$t=3s,\eta=u^2/9$を代入し、$x,y,z$を$\alpha,\beta,\gamma$に置き換える。
中心スカラー$2u^3/9$を加えると表示式を得る。
\end{proof}

\subsection{相対ポテンシャル複体の実際の完全性}
$Q=Q_\epsilon$、$S=Q[u]$とし、$D=D_\epsilon$を$S$-代数とみる。
自由$S$-代数のcyclic potentialを
\begin{equation}
 W=xyz-yxz-h\operatorname{sym}(f_\epsilon+ux^2),
 \qquad r_i=\partial^{\mathrm{cyc}}_{x_i}W,\quad(x_1,x_2,x_3)=(x,y,z)
\end{equation}
とする。定義関係式は$r_1=r_2=r_3=0$である。

$\partial^{\mathrm{nc}}_j$は語の各$x_j$を取り除き、その左右の部分を分ける演算とする：
\begin{equation}
 \partial^{\mathrm{nc}}_j(a_1\cdots a_d)
 =\sum_{a_\ell=x_j} a_1\cdots a_{\ell-1}\otimes a_{\ell+1}\cdots a_d.
\end{equation}
係数環$S$は微分しない。
\begin{proposition}[底変更と両立する第一ブロックの相対対角]
\label{nq:prop:relative-diagonal}
$D^e_S=D\otimes_S D^{\mathrm{op}}$とする。有限自由双加群の複体
\begin{equation}
 0\longrightarrow D^e_S
 \xrightarrow{d_3}(D^e_S)^3
 \xrightarrow{d_2}(D^e_S)^3
 \xrightarrow{d_1}D^e_S
 \xrightarrow{m}D\longrightarrow0
 \label{nq:eq:potential-resolution}
\end{equation}
を、基底$e_i$（生成元）、$\varepsilon_i$（関係式）について
\begin{equation}\begin{aligned}
 d_1(e_i)&=x_i\otimes1-1\otimes x_i,\\
 d_2(\varepsilon_i)&=\sum_j\sum_w\ \sum_{w=p x_j q} c_{iw}p e_j q,\\
 d_3(1)&=\sum_i(x_i\varepsilon_i-\varepsilon_i x_i)
\end{aligned}\end{equation}
で定める。ここで$r_i=\sum_w c_{iw}w$と展開している。
この複体は完全で、任意の可換$S$-代数への底変更後も完全である。
\end{proposition}
\begin{proof}
$d_1d_2=0$は$r_i=0$から、$d_2d_3=0$はcyclic potentialの恒等式
$\sum_i[x_i,r_i]=0$を非可換微分することから従う。
active次数に関して、複体の自由基底のfiltration次数を$0,1,2,3$と取る。
$ux^2$はactive次数二なので、随伴次数付き複体は$u=0$の同次ポテンシャルの
Koszul双加群分解を$Q[u]$へ拡大したものである。

この同次複体が$Q$上で完全であることを確認する。
補題\ref{nq:lem:pbw}の$u=0$で、各次数成分は有限自由$Q$-加群である。
各複素点に特殊化すると、Mori--Smithの三次元Calabi--Yau代数に対する
同次ポテンシャル分解となり完全である
\cite[Theorem 4.2]{MSNE12}\cite[\S3.2]{BTNE12}。
各次数における複体は有限個の有限自由$Q$-加群からなり、
全閉点のfiberで完全なので$Q$上で完全である。
さらに最後の項も自由であるから、各次数の完全列は$Q$-加群として分裂する。
従って任意の底変更後にも完全性が保たれる。

filtrationは下に有界で、各元のactive次数は有限である。
随伴複体でcycleの最高次数を消し、次数を下げる帰納法により
元の複体の完全性を得る。底変更後も同じmonic PBW表示と同じ随伴分解を持つため、
同じ証明が適用できる。これは既知のPBWポテンシャル分解
\cite[Theorem 3.6の証明]{BTNE12}を、今回の係数環と底変更について明示したものである。
\end{proof}

\begin{remark}[中心方向の通常微分は使えない]\label{nq:rem:no-derivation}
$h\ne0$のファイバーでは、$\partial_u(u)=1$を満たす次数$-1$の
代数導分$\partial_u:D_\epsilon\to D_\epsilon$は存在しない。
実際、次数から$\partial_u(x),\partial_u(y),\partial_u(z)$はスカラーとなる。
尖点型の$[y,z]-h(z^2+2ux)=0$を微分して$u=0$に落とすと、
$-2h(x+\partial_u(z)z)=0$という一次の関係が必要になるが、PBW性に反する。
節点型でも$-h(2x+\partial_u(y)z+\partial_u(z)y)=0$となって同様である。
相対対角\eqref{nq:eq:potential-resolution}は$u$を係数としており、
この不存在とは矛盾しない。中心三次商の残余圏を調べる際に、
可換多項式の$\partial_u$をそのまま作用させることはできない。
\end{remark}

\subsection{七生成元への組立てと中心三次商}
まず$E_s$を$\CC[s]$上で考える。$u=0$のFermat型三次元Jacobian代数は
全ての$s$で三次元Calabi--Yau AS-regularである
\cite[Theorem 4.2(1)]{MSNE12}。
$u$を加えた関係式の一次混合係数は対角行列であり、
cyclic potentialのsyzygyにより
\cite[Theorem 3.1.3]{LBSVNE12}の条件を$\gamma=0$で満たす。
従って各$s$で$E_s$は中心正則$u$を持つ四次元Koszul AS-regular domainである。
各次数の$\CC[s]$-加群は有限生成で全fiberの次元が一定なので自由となる。
$u$による完全列を次数ごとに分裂させると、$E_s$は$\CC[s,u]$上自由で、
active filtrationの随伴次数付き環は同次Fermat代数の係数拡大となる。
そのポテンシャル複体も命題\ref{nq:prop:relative-diagonal}と同じ議論で、
任意の底変更と両立する長さ三の相対自由対角分解を与える。

\begin{proof}[定理\ref{nq:thm:family}の証明]
中心三次元を
\begin{equation}
 \Omega_\epsilon=\Gamma_\epsilon(r)+\Gamma_{\mathrm F}(s)+u^3
 \label{nq:eq:total-center}
\end{equation}
と定める。中心性と原点への特殊化は既に示した。
両ブロックが$R_\epsilon[u]$上自由なので、そのテンソル積も自由である。
全体の$u=0$での商は、二つの三次元Koszul AS-regular代数のテンソル積であり、
Hilbert級数$(1-t)^{-6}$、AS型$(6,6)$を持つ。
このテンソル積がnoetherian domainであることには注意が要る。
一般のnoetherian代数のテンソル積についての無条件な主張は用いず、
第一ブロックの上記Ore表示を第二のFermat代数上へ係数拡大する。
これにより二変数量子平面またはJordan平面と、もう一つの可逆Ore拡大の反復として得られる。

$u$は全体でも中心正則元である。
中心正則次数一拡大に対するKoszul性の持上げ
\cite[Theorem 2.6, Remark 2.8]{LBSVNE12}と、
$H_{\mathcal A}(t)=(1-t)^{-7}$から、Koszul双対の級数は$(1+t)^7$となりglobal dimensionは七である。
AS-Gorenstein型は中心正則元に対するchange-of-ringsの
\begin{equation}
 \operatorname{Ext}_{\mathcal A}^{i}(k,\mathcal A)
 \simeq\operatorname{Ext}_{\mathcal A/(u)}^{i-1}
             (k,\mathcal A/(u))(1)
\end{equation}
および左側の同じ式から$(7,7)$となる。
noetherian性と域性は$u$-進次数に関する随伴環
$(\mathcal A/(u))[u]$から従う。正次数の$u$を用いるので、固定した全次数における
最高項の消去は有限回で停止する。

各ファイバーで$\Omega_\epsilon\bmod u$の第一ブロック次数$(3,0)$の項は
$\Gamma_\epsilon(r)|_{u=0}$であり、そのPBW正規形は非零である。
第二ブロック次数$(0,3)$の項とは相殺できないため、$\Omega_\epsilon\ne0$である。
周囲がdomainなので左右正則である。
各次数で$\Omega_\epsilon:\mathcal A_{d-3}\to\mathcal A_d$は
全ての幾何学的fiberで単射だから、局所的に分裂した単射であり、
余核は有限局所自由で任意の底変更と両立する。
中心正則次数三商のchange-of-ringsからAS型$(6,4)$を得る。
半古典極限は関係式を$h$で割って$r=h\lambda,s=h\mu$の原点へ制限すればよい。

最後に、$R_\epsilon[u]$上で二つの長さ三の相対ポテンシャル分解をテンソルしtotalizeする。
その自由階数列は$(1+t)^3(1+t)^3=(1+t)^6$であり、完全性と任意の底変更が保たれる。
$R_\epsilon$上では包絡代数から相対包絡代数への商は中心元
$u\otimes1-1\otimes u$による商であり、この元は正則である。
その長さ一のKoszul分解と相対分解から、周囲の対角は絶対的にもperfect（射影次元高々七）となる。
\end{proof}

\subsection{相対残余圏について残る仕事}
本節は二つの固定Poisson成分の稠密開を同時に扱う量子化族を与える。
一つの方向だけを固定した形式弧ではなく、$(r,s)$の両方向を保持する。
周囲の代数の有限対角分解は実在し、その底変更も証明した。
しかし、中心三次商のaffine coneは可換点でも原点で特異である。
\emph{周囲の対角を中心三次商でテンソルするだけ}では、その商や残余圏の対角にはならない。
次に必要なのは、中心三次商の相対局所cohomology・射影化とOrlov windowを含むproperなdg模型、
およびその模型の対角のperfect性である。
この過程で中心方向の混合を除かず、注意\ref{nq:rem:no-derivation}の障害を扱わなければならない。
一般次元・全Poisson成分への可換復元は、引き続きこの未解決段階を含む。